\documentclass[11pt]{article}

\usepackage{amsmath,amssymb}
\usepackage{xurl}
\usepackage[colorlinks=true,linkcolor=blue,citecolor=blue,urlcolor=blue]{hyperref}

% Shared commands for the notes in docs/paper-gaps/.
% Every note loads this file via % Shared commands for the notes in docs/paper-gaps/.
% Every note loads this file via % Shared commands for the notes in docs/paper-gaps/.
% Every note loads this file via \input{command}. Project-wide notation
% belongs here, never in per-note redefinitions.

% --- Project configuration ------------------------------------------------
\newcommand{\projectIssueBase}{https://github.com/LionSR/QICLean/issues/}
\newcommand{\projectPrBase}{https://github.com/LionSR/QICLean/pull/}
\newcommand{\projectDocBase}{https://sirui-lu.com/QICLean/docs/}
% ---------------------------------------------------------------------------

\DeclareUnicodeCharacter{2080}{\ifmmode{}_{0}\else\textsubscript{0}\fi}
\DeclareUnicodeCharacter{2081}{\ifmmode{}_{1}\else\textsubscript{1}\fi}
\DeclareUnicodeCharacter{2082}{\ifmmode{}_{2}\else\textsubscript{2}\fi}
\DeclareUnicodeCharacter{2099}{\ifmmode{}_{n}\else\textsubscript{n}\fi}

\newcommand{\Lean}{\textsc{Lean}}

% A formal declaration name, upright sans, breakable at underscores.
\ExplSyntaxOn
\NewDocumentCommand{\leanid}{m}{
  \ifmmode
    \textsf{\detokenize{#1}}
  \else
    \group_begin:
    \urlstyle{sf}
    \tl_set:Nn \l_tmpa_tl {#1}
    \tl_replace_all:Nnn \l_tmpa_tl {\_} {_}
    \exp_args:NV \path \l_tmpa_tl
    \group_end:
  \fi
}
\ExplSyntaxOff

% Traceability links; use in footnotes, not in the running prose.
\newcommand{\ghissue}[1]{\href{\projectIssueBase#1}{GitHub issue \##1}}
\newcommand{\ghpr}[1]{\href{\projectPrBase#1}{PR \##1}}
\newcommand{\doclink}[1]{\href{\projectDocBase#1.html}{\path{#1}}}

% Domain notation, using \providecommand so a note-local refinement stays
% possible.
\providecommand{\tr}{\operatorname{tr}}

% Dirac notation and the identity operator, as in blueprint/src/macros.
\usepackage{dsfont}
\providecommand{\ket}[1]{|#1\rangle}
\providecommand{\bra}[1]{\langle #1|}
\providecommand{\braket}[2]{\langle #1 | #2 \rangle}
\providecommand{\ketbra}[2]{|#1\rangle\!\langle #2|}
\providecommand{\Id}{\mathds{1}}
\providecommand{\one}{\mathds{1}}
\providecommand{\C}{\mathbb{C}}
\providecommand{\MN}[1]{M_{#1}(\C)}
\providecommand{\MD}{M_D(\C)}

% Verdict marker: \gapnote{<kind>}{<status>}. Kind is the policy
% classification (clarification, local-correction, scope-restriction,
% unfaithful, false-source, open-gap); status is open, wip (elimination
% underway), resolved, or historical. Machine-read by the site index;
% prints a verdict line.
\newcommand{\gapnote}[2]{\par\noindent\textbf{Verdict:} #1 (#2).\par}
. Project-wide notation
% belongs here, never in per-note redefinitions.

% --- Project configuration ------------------------------------------------
\newcommand{\projectIssueBase}{https://github.com/LionSR/QICLean/issues/}
\newcommand{\projectPrBase}{https://github.com/LionSR/QICLean/pull/}
\newcommand{\projectDocBase}{https://sirui-lu.com/QICLean/docs/}
% ---------------------------------------------------------------------------

\DeclareUnicodeCharacter{2080}{\ifmmode{}_{0}\else\textsubscript{0}\fi}
\DeclareUnicodeCharacter{2081}{\ifmmode{}_{1}\else\textsubscript{1}\fi}
\DeclareUnicodeCharacter{2082}{\ifmmode{}_{2}\else\textsubscript{2}\fi}
\DeclareUnicodeCharacter{2099}{\ifmmode{}_{n}\else\textsubscript{n}\fi}

\newcommand{\Lean}{\textsc{Lean}}

% A formal declaration name, upright sans, breakable at underscores.
\ExplSyntaxOn
\NewDocumentCommand{\leanid}{m}{
  \ifmmode
    \textsf{\detokenize{#1}}
  \else
    \group_begin:
    \urlstyle{sf}
    \tl_set:Nn \l_tmpa_tl {#1}
    \tl_replace_all:Nnn \l_tmpa_tl {\_} {_}
    \exp_args:NV \path \l_tmpa_tl
    \group_end:
  \fi
}
\ExplSyntaxOff

% Traceability links; use in footnotes, not in the running prose.
\newcommand{\ghissue}[1]{\href{\projectIssueBase#1}{GitHub issue \##1}}
\newcommand{\ghpr}[1]{\href{\projectPrBase#1}{PR \##1}}
\newcommand{\doclink}[1]{\href{\projectDocBase#1.html}{\path{#1}}}

% Domain notation, using \providecommand so a note-local refinement stays
% possible.
\providecommand{\tr}{\operatorname{tr}}

% Dirac notation and the identity operator, as in blueprint/src/macros.
\usepackage{dsfont}
\providecommand{\ket}[1]{|#1\rangle}
\providecommand{\bra}[1]{\langle #1|}
\providecommand{\braket}[2]{\langle #1 | #2 \rangle}
\providecommand{\ketbra}[2]{|#1\rangle\!\langle #2|}
\providecommand{\Id}{\mathds{1}}
\providecommand{\one}{\mathds{1}}
\providecommand{\C}{\mathbb{C}}
\providecommand{\MN}[1]{M_{#1}(\C)}
\providecommand{\MD}{M_D(\C)}

% Verdict marker: \gapnote{<kind>}{<status>}. Kind is the policy
% classification (clarification, local-correction, scope-restriction,
% unfaithful, false-source, open-gap); status is open, wip (elimination
% underway), resolved, or historical. Machine-read by the site index;
% prints a verdict line.
\newcommand{\gapnote}[2]{\par\noindent\textbf{Verdict:} #1 (#2).\par}
. Project-wide notation
% belongs here, never in per-note redefinitions.

% --- Project configuration ------------------------------------------------
\newcommand{\projectIssueBase}{https://github.com/LionSR/QICLean/issues/}
\newcommand{\projectPrBase}{https://github.com/LionSR/QICLean/pull/}
\newcommand{\projectDocBase}{https://sirui-lu.com/QICLean/docs/}
% ---------------------------------------------------------------------------

\DeclareUnicodeCharacter{2080}{\ifmmode{}_{0}\else\textsubscript{0}\fi}
\DeclareUnicodeCharacter{2081}{\ifmmode{}_{1}\else\textsubscript{1}\fi}
\DeclareUnicodeCharacter{2082}{\ifmmode{}_{2}\else\textsubscript{2}\fi}
\DeclareUnicodeCharacter{2099}{\ifmmode{}_{n}\else\textsubscript{n}\fi}

\newcommand{\Lean}{\textsc{Lean}}

% A formal declaration name, upright sans, breakable at underscores.
\ExplSyntaxOn
\NewDocumentCommand{\leanid}{m}{
  \ifmmode
    \textsf{\detokenize{#1}}
  \else
    \group_begin:
    \urlstyle{sf}
    \tl_set:Nn \l_tmpa_tl {#1}
    \tl_replace_all:Nnn \l_tmpa_tl {\_} {_}
    \exp_args:NV \path \l_tmpa_tl
    \group_end:
  \fi
}
\ExplSyntaxOff

% Traceability links; use in footnotes, not in the running prose.
\newcommand{\ghissue}[1]{\href{\projectIssueBase#1}{GitHub issue \##1}}
\newcommand{\ghpr}[1]{\href{\projectPrBase#1}{PR \##1}}
\newcommand{\doclink}[1]{\href{\projectDocBase#1.html}{\path{#1}}}

% Domain notation, using \providecommand so a note-local refinement stays
% possible.
\providecommand{\tr}{\operatorname{tr}}

% Dirac notation and the identity operator, as in blueprint/src/macros.
\usepackage{dsfont}
\providecommand{\ket}[1]{|#1\rangle}
\providecommand{\bra}[1]{\langle #1|}
\providecommand{\braket}[2]{\langle #1 | #2 \rangle}
\providecommand{\ketbra}[2]{|#1\rangle\!\langle #2|}
\providecommand{\Id}{\mathds{1}}
\providecommand{\one}{\mathds{1}}
\providecommand{\C}{\mathbb{C}}
\providecommand{\MN}[1]{M_{#1}(\C)}
\providecommand{\MD}{M_D(\C)}

% Verdict marker: \gapnote{<kind>}{<status>}. Kind is the policy
% classification (clarification, local-correction, scope-restriction,
% unfaithful, false-source, open-gap); status is open, wip (elimination
% underway), resolved, or historical. Machine-read by the site index;
% prints a verdict line.
\newcommand{\gapnote}[2]{\par\noindent\textbf{Verdict:} #1 (#2).\par}


\title{Cross-Band Commutation in the Transport Estimate of the Area-Law Paper}
\author{}
\date{2026-10-08}

\begin{document}
\maketitle
\gapnote{scope-restriction}{open}

\section{The source hypothesis}

Proposition~7.4 of Ref.~\cite{OpenAI2026AreaLaw} (\path{06-transport.tex},
lines~377--434) concerns band metrics
$A(P,Y,F)=(W_P^{-1}W_F^{-1}W_Y)^2$ on the symmetric subspace
$\mathcal S_k=\mathrm{Sym}^k\mathcal V$ (lines~251--258). Its commutation
condition (lines~273--276) asks that every old or new metric of band $g$
commute with every old or new metric of band $g'\neq g$. Since the metrics
are operators on $\mathcal S_k$, this is commutation on $\mathcal S_k$. The
paper uses this point when it applies the proposition: under the nesting
condition of \path{07-comparators.tex} (lines~26--35), the different bands'
metric algebras commute only after complementary labels are identified on the
symmetric space, using $W_{Q,k}=W_{Q^c,k}$ on $\mathcal S_k$
(\path{05-replicas.tex}, line~312).

\section{The formal restriction}

The formal band metric is the operator
$(W_P^{-1}W_F^{-1}W_Y)^2$ on the full tensor power
$\mathcal V^{\otimes k}$. It commutes with the copy permutations, so it
preserves $\mathcal S_k$ and restricts there to the source metric. The formal
cross-band hypothesis asks for commutation of these operators on all of
$\mathcal V^{\otimes k}$, for every $k$.\footnote{Formal declarations:
\leanid{TensorPower.ReplicaTransport.TransportData.CrossBandCommute} in
\path{QICLean/Representation/ReplicaTransport/Setup.lean}, assumed for every $k$ by
\leanid{TensorPower.ReplicaTransport.transport_of_relativePin_of_splitSkewBound}.}

Commutation on the full tensor power implies commutation on $\mathcal S_k$,
but the converse fails. Central label functions of two subsystems $Q,Q'$
commute on $\mathcal V^{\otimes k}$ when $Q\subseteq Q'$ or $Q\cap Q'=\emptyset$,
but in general not when the two subsystems overlap without nesting. For three
qubits $a,b,c$, $k=3$, $t=1/10$ and the label weights of
\path{05-replicas.tex} (display \texttt{replicas:w-definition}, padding
dimension $8$), a direct computation gives
$\|[W_{ab},W_{bc}]\|_2\approx 3.9\times10^{-4}$ on $\mathcal V^{\otimes 3}$,
while the commutator vanishes on $\mathcal S_3$, where
$W_{ab}=W_c$ and $W_{bc}=W_a$.

In the comparator setting, $P_g=CXU_g$, $F_g=RV_g$, and for $g<g'$ the nesting
condition gives $XU_gY_g\subseteq XU_{g'}$. Every pair of factors from the two
bands is nested or disjoint, except $F_g=RV_g$ and $P_{g'}=CXU_{g'}$, which
overlap in $V_g\cap U_{g'}$ without nesting. The remaining factors of each
band commute with every factor of the other band, so the full-space
commutation of $A_g$ and $A_{g'}$ reduces to that of $W_{RV_g}$ and
$W_{CXU_{g'}}$, which fails in general. The formal theorem therefore proves
Proposition~7.4 under a commutation hypothesis that the source does not
impose and that its multi-band application need not satisfy; for one band,
or for bands whose factors are pairwise nested or disjoint, the two
hypotheses agree in effect.

\section{Elimination plan}

Replace each band metric in the tree by
$\Pi A\Pi+(I-\Pi)$, where $\Pi$ is the projection onto $\mathcal S_k$. This
operator is positive definite, agrees with $A$ on $\mathcal S_k$, and two such
operators commute on $\mathcal V^{\otimes k}$ exactly when the restrictions of
the original metrics to $\mathcal S_k$ commute. Restating the cross-band
hypothesis for these operators gives the source hypothesis. The positivity,
logarithm factorization and energy arguments then go through unchanged on
$\mathcal S_k$, and the relative pin and skew bound are needed only on
symmetric vectors and symmetric densities, which is how they are already
stated.

\section{Verdict}

\textbf{Verdict: scope restriction.} The conditional transport theorem proves
the conclusions of Proposition~7.4 of Ref.~\cite{OpenAI2026AreaLaw} assuming
cross-band commutation on the full tensor power rather than on the symmetric
subspace. The source-faithful statement remains open.

\bibliographystyle{alpha}
\bibliography{references}

\end{document}
